\chapter*{第 8 章 黏性流体动力学基础}\addcontentsline{toc}{chapter}{第 8 章 黏性流体动力学基础}\markboth{第 8 章 黏性流体动力学基础 解答}{}

\anstitle{第 8 章 习题解答}

\sloppy
\emergencystretch=2em

\begin{tishi}
\textbf{题源说明：}本章收录的是\textbf{孔珑《工程流体力学》（第四版）第 9 章"黏性流体绕过物体的流动"
的课后习题中符合 818 考纲的题目全套解答}（孔珑原书题号 9-9、9-10、9-11、9-12、9-14、9-15、
9-16、9-17、9-18、9-19、9-20，共 11 题），按本册归章规则列入\textbf{赵琴教材第 8 章 黏性流体动力学基础}。
解答依据《工程流体力学 学习指导及习题解答》（陈洁、袁铁江 编著）第 9 章"课后习题解答"，
并按本册体例补入"已知—依据—步骤—结果—易错提醒"五段。
全书数据与原书一致；原书有两处数值自相矛盾（第 8.2 题的雷诺数、第 8.6 题的板末端雷诺数），
已在相应题的步骤中给出\textbf{两种口径}，并在题目册章末集中说明。
本章略去的题号（9-1 至 9-8 属缝隙层流，归赵琴第 5 章；9-13 原书为"略"；9-21 属自由射流）
及其理由，见题目册。
\end{tishi}

\noindent\textbf{第 8.1 题（孔珑 9-9）\quad 求空气纵向绕流平板时边界层的厚度}

\begin{jie}
\textbf{已知：}平板极薄、沿来流方向放置，纵向绕流；来流速度 $v_\infty=30\ \mathrm{m/s}$；
空气温度 $25\ ^{\circ}\mathrm{C}$，运动黏度 $\nu=16\times10^{-6}\ \mathrm{m^2/s}$；
量取位置距平板前缘 $x=200\ \mathrm{mm}=0.2\ \mathrm{m}$；压强为大气压（不影响运动黏度取值）。

\textbf{求：}$x=200\ \mathrm{mm}$ 处的边界层厚度 $\delta$。

\textbf{依据：}先由 $\Rey_x=v_\infty x/\nu$ 判断流态，平板边界层由层流转为湍流的临界雷诺数
$\Rey_{x\mathrm{cr}}=5\times10^{5}$；层流边界层厚度按布拉修斯解（教材 8.4 节）
\[ \frac{\delta}{x}=\frac{5.84}{\sqrt{\Rey_x}},\qquad \delta=\frac{5.84x}{\sqrt{\Rey_x}} \]

\textbf{第 1 步\quad 计算局部雷诺数}
\[ \Rey_x=\frac{v_\infty x}{\nu}=\frac{30\times0.2}{16\times10^{-6}}=\frac{6}{16\times10^{-6}}=3.75\times10^{5} \]
因为 $3.75\times10^{5}<5\times10^{5}$，故 $x=200\ \mathrm{mm}$ 处仍为\textbf{层流边界层}。

\textbf{第 2 步\quad 代入层流边界层厚度公式}
\[ \delta=\frac{5.84\times0.2}{\sqrt{3.75\times10^{5}}}=\frac{1.168}{612.4}=1.907\times10^{-3}\ \mathrm{m}=1.91\ \mathrm{mm} \]

\textbf{结果：}离平板前缘 $200\ \mathrm{mm}$ 处边界层厚度 $\delta=1.907\ \mathrm{mm}\approx1.91\ \mathrm{mm}$。

\textbf{量纲自检：}$\Rey_x$ 为无量纲数；$\delta$ 的分子为长度（$\mathrm{m}$），分母无量纲，
故单位为 $\mathrm{m}$，量纲正确 $\checkmark$。
\end{jie}

\begin{yicuo}
\textbf{易错点一：运动黏度必须与温度对应。}
$25\ ^{\circ}\mathrm{C}$ 空气取 $\nu=16\times10^{-6}\ \mathrm{m^2/s}$；
若误取 $20\ ^{\circ}\mathrm{C}$ 的 $15\times10^{-6}\ \mathrm{m^2/s}$，$\delta$ 会偏大 3.3%，
至于把液体的黏度拿来用则完全错。

\textbf{易错点二：边界层厚度与 $\sqrt{x}$ 成正比，不是与 $x$ 成正比。}
由 $\delta=5.84\sqrt{\nu x/v_\infty}$ 可见 $\delta\propto\sqrt{x}$；
"离前缘越远越厚"是对的，但按线性增长估算会明显偏大。

\textbf{易错点三：$5.84$ 这个系数有"口径"问题。}
它是布拉修斯精确解的结果；用假定速度剖面（抛物线 $2\eta-\eta^{2}$、正弦 $\sin(\pi\eta/2)$）
代入动量积分方程分别得到 $5.48$ 和 $4.79$。
2021 年真题卷面给出的是 $5.48/\sqrt{\Rey_x}$，\textbf{卷面给哪个就用哪个}。
\end{yicuo}

\noindent\textbf{第 8.2 题（孔珑 9-10）\quad 油流纵向绕流薄平板的总摩擦阻力与边界层厚度}

\begin{jie}
\textbf{已知：}油温 $20\ ^{\circ}\mathrm{C}$，密度 $\rho=925\ \mathrm{kg/m^3}$，
运动黏度 $\nu=7.9\times10^{-5}\ \mathrm{m^2/s}$；来流速度 $v_\infty=60\ \mathrm{cm/s}=0.6\ \mathrm{m/s}$；
平板长 $L=50\ \mathrm{cm}=0.5\ \mathrm{m}$、宽 $b=15\ \mathrm{cm}=0.15\ \mathrm{m}$。

\textbf{求：}平板总摩擦阻力 $F$（两面受油流剪切）和板末端边界层厚度 $\delta$。

\textbf{依据：}由 $\Rey_l=v_\infty L/\nu$ 判断流态；层流平板单面摩擦阻力
\[ F_{D\mathrm{f}}=C_{\mathrm{Df}}\cdot\frac{1}{2}\rho v_\infty^{2}\cdot bL,\qquad
C_{\mathrm{Df}}=\frac{1.328}{\sqrt{\Rey_l}}\ \Longrightarrow\ F_{D\mathrm{f}}=0.664\,bL\rho v_\infty^{2}\Rey_l^{-1/2} \]
原书代入时把系数写成 $0.686$（相当于取 $C_{\mathrm{Df}}=1.372/\sqrt{\Rey_l}$），本册按原书口径计算并同时给出标准口径；
板末端边界层厚度按层流公式 $\delta=5.84L/\sqrt{\Rey_l}$。

\textbf{第 1 步\quad 计算板末端雷诺数}
\[ \Rey_l=\frac{v_\infty L}{\nu}=\frac{0.6\times0.5}{7.9\times10^{-5}}=\frac{0.3}{7.9\times10^{-5}}=3.797\times10^{3}<5\times10^{5} \]
故全板为\textbf{层流边界层}。
（原书在求厚度处把该雷诺数印作 $3.797\times10^{4}$，与本题的 $10^{3}$ 自相矛盾；
原书给出的 $\delta=47.387\ \mathrm{mm}$ 正是用 $3.797\times10^{3}$ 算出来的，故正确值为 $3.797\times10^{3}$。）

\textbf{第 2 步\quad 平板一面上的摩擦阻力（原书口径，系数 $0.686$）}
\[ F_{D\mathrm{f}}=0.686\,bL\rho v_\infty^{2}\Rey_l^{-1/2}
=0.686\times0.15\times0.5\times925\times0.6^{2}\times(3.797\times10^{3})^{-1/2} \]
\[ =0.686\times0.075\times925\times0.36\times0.01623=0.278\ \mathrm{N} \]

\textbf{第 3 步\quad 平板两面的总摩擦阻力}
\[ F=2F_{D\mathrm{f}}=2\times0.278=0.556\ \mathrm{N} \]

\textbf{第 4 步\quad 板末端边界层厚度}
\[ \delta=\frac{5.84L}{\sqrt{\Rey_l}}=\frac{5.84\times0.5}{\sqrt{3.797\times10^{3}}}
=\frac{2.92}{61.62}=0.04739\ \mathrm{m}=47.39\ \mathrm{mm} \]

\textbf{结果：}一面摩擦阻力 $0.278\ \mathrm{N}$，平板两面总摩擦阻力 $F=0.556\ \mathrm{N}$；
板末端边界层厚度 $\delta=47.387\ \mathrm{mm}\approx47.4\ \mathrm{mm}$。
若按教材通行的 $C_{\mathrm{Df}}=1.328/\sqrt{\Rey_l}$（系数 $0.664$）计算，
则一面 $0.269\ \mathrm{N}$、两面 $0.538\ \mathrm{N}$，比原书口径小 3.3%。

\textbf{量纲自检：}$F=0.686\,bL\rho v_\infty^{2}=[\mathrm{m}][\mathrm{m}][\mathrm{kg/m^3}][\mathrm{m^2/s^2}]
=\mathrm{kg\cdot m/s^2}=\mathrm{N}$，量纲正确 $\checkmark$；$\delta$ 为 $\mathrm{m}$ $\checkmark$。
\end{jie}

\begin{yicuo}
\textbf{易错点一：$0.686$ 与 $0.664$ 两个系数来源不同。}
原书取 $0.686$（对应 $C_{\mathrm{Df}}=1.372/\sqrt{\Rey_l}$），
教材与 2021 年真题卷面给的是 $C_{\mathrm{Df}}=1.328/\sqrt{\Rey_l}$（对应 $0.664$）。
两者只差 3.3%，但答题时应\textbf{以题给公式为准}，不要两个系数混用。

\textbf{易错点二：题目问"总摩擦阻力"，必须乘 2。}
薄平板两面都与油接触，油流在两面都产生摩擦阻力，原书答案 $0.556\ \mathrm{N}$ 就是两面的值；
只算一面会丢掉一半。

\textbf{易错点三：边界层厚度要用板长算的 $\Rey_l$。}
厚度按 $x=L$ 处的 $\Rey_l$ 计算，不能用某个局部 $x$ 的 $\Rey_x$；
本题 $\Rey_l$ 只有 $3.8\times10^{3}$，所以 $0.5\ \mathrm{m}$ 长的板上边界层已厚达 $47\ \mathrm{mm}$，
这也说明\textbf{低雷诺数下层流边界层可能相当厚}，不能凭"边界层很薄"的印象去估算。
\end{yicuo}

\noindent\textbf{第 8.3 题（孔珑 9-11）\quad 抛物线速度分布下层流边界层的厚度与摩擦阻力系数}

\begin{jie}
\textbf{已知：}平板层流边界层内速度分布 $v_x/v_\infty=2\,(y/\delta)-(y/\delta)^{2}$（抛物线分布）；
来流速度 $v_\infty$、运动黏度 $\nu$；平板长 $L$、宽 $b$。

\textbf{求：}边界层厚度 $\delta$、摩擦阻力系数 $C_{\mathrm{Df}}$ 与雷诺数 $\Rey$ 的关系式。

\textbf{依据：}
① 壁面切应力由牛顿内摩擦定律 $\tau_w=\mu\left(\partial v_x/\partial y\right)_{y=0}$ 求；
② 冯·卡门动量积分方程
\[ \frac{\mathrm{d}\delta_2}{\mathrm{d}x}=\frac{\tau_w}{\rho v_\infty^{2}},\qquad
\delta_2=\int_0^{\delta}\frac{v_x}{v_\infty}\left(1-\frac{v_x}{v_\infty}\right)\mathrm{d}y \]
其中 $\delta_2$ 为动量损失厚度（教材 8.5 节）。

\textbf{第 1 步\quad 壁面切应力}
令 $\eta=y/\delta$，则 $v_x=v_\infty\left[2\eta-\eta^{2}\right]$，
\[ \frac{\partial v_x}{\partial y}=v_\infty\left[\frac{2}{\delta}-\frac{2y}{\delta^{2}}\right] \]
\[ \tau_w=\mu\left(\frac{\partial v_x}{\partial y}\right)_{y=0}=\mu v_\infty\cdot\frac{2}{\delta}=\frac{2\mu v_\infty}{\delta} \qquad (8\text{-}1) \]

\textbf{第 2 步\quad 动量损失厚度}
\[ \delta_2=\delta\int_0^{1}\left(2\eta-\eta^{2}\right)\left[1-\left(2\eta-\eta^{2}\right)\right]\mathrm{d}\eta
=\delta\int_0^{1}\left(2\eta-\eta^{2}\right)(1-\eta)^{2}\mathrm{d}\eta \]
\[ \left(2\eta-\eta^{2}\right)(1-\eta)^{2}=2\eta-5\eta^{2}+4\eta^{3}-\eta^{4} \]
\[ \delta_2=\delta\left[1-\frac{5}{3}+1-\frac{1}{5}\right]=\frac{2}{15}\delta=0.1333\delta \]

\textbf{第 3 步\quad 代入动量积分方程}
\[ \frac{2}{15}\frac{\mathrm{d}\delta}{\mathrm{d}x}=\frac{\tau_w}{\rho v_\infty^{2}}
=\frac{2\mu v_\infty}{\delta\rho v_\infty^{2}}=\frac{2\nu}{v_\infty\delta} \]
化简得
\[ \delta\,\mathrm{d}\delta=15\frac{\nu}{v_\infty}\mathrm{d}x \]

\textbf{第 4 步\quad 积分并定积分常数}
\[ \frac{1}{2}\delta^{2}=15\frac{\nu}{v_\infty}x+C \]
边界条件：$x=0$ 处 $\delta=0$，故 $C=0$。
\[ \frac{1}{2}\delta^{2}=15\frac{\nu}{v_\infty}x
\ \Longrightarrow\ \delta=\sqrt{30}\sqrt{\frac{\nu x}{v_\infty}}
=5.48\sqrt{\frac{\nu x}{v_\infty}} \]
\[ \delta=\frac{5.48x}{\sqrt{\Rey_x}} \]

\textbf{第 5 步\quad 壁面切应力与雷诺数的关系}
\[ \tau_w=\frac{2\mu v_\infty}{\delta}=\frac{2\mu v_\infty}{5.48}\sqrt{\frac{v_\infty}{\nu x}}
=0.365\,\mu v_\infty^{3/2}\nu^{-1/2}x^{-1/2} \]
利用 $\mu=\rho\nu$：
\[ \tau_w=0.365\,\rho v_\infty^{2}\left(\frac{\nu}{v_\infty x}\right)^{1/2}
=0.365\,\rho v_\infty^{2}\Rey_x^{-1/2} \]

\textbf{第 6 步\quad 长 $L$、宽 $b$ 的平板一面的总摩擦阻力}
\[ F_{D\mathrm{f}}=b\int_0^{L}\tau_w\mathrm{d}x
=0.365b\rho v_\infty^{2}\sqrt{\frac{\nu}{v_\infty}}\int_0^{L}x^{-1/2}\mathrm{d}x
=0.365b\rho v_\infty^{2}\sqrt{\frac{\nu}{v_\infty}}\cdot2\sqrt{L} \]
\[ F_{D\mathrm{f}}=0.73\,bL\rho v_\infty^{2}\Rey_l^{-1/2} \]

\textbf{第 7 步\quad 摩擦阻力系数（分母用平板一面面积 $bL$ 与动压）}
\[ C_{\mathrm{Df}}=\frac{F_{D\mathrm{f}}}{\frac{1}{2}\rho v_\infty^{2}bL}
=\frac{0.73\times2}{\sqrt{\Rey_l}}=\frac{1.46}{\sqrt{\Rey_l}} \]

\textbf{结果：}$\delta=\dfrac{5.48x}{\sqrt{\Rey_x}}$；
$\tau_w=0.365\rho v_\infty^{2}\Rey_x^{-1/2}$；
$F_{D\mathrm{f}}=0.73\,bL\rho v_\infty^{2}\Rey_l^{-1/2}$；
$C_{\mathrm{Df}}=\dfrac{1.46}{\sqrt{\Rey_l}}$。

\textbf{量纲自检：}$\delta_2$ 与 $\delta$ 同为 $\mathrm{m}$；
$\tau_w/(\rho v_\infty^{2})=[\mathrm{Pa}]/([\mathrm{kg/m^3}][\mathrm{m^2/s^2}])=\mathrm{m}$，
与 $\mathrm{d}\delta_2/\mathrm{d}x$ 的量纲一致 $\checkmark$；$C_{\mathrm{Df}}$ 为无量纲 $\checkmark$。
\end{jie}

\begin{yicuo}
\textbf{易错点一：动量损失厚度 $\delta_2$ 不要记错。}
抛物线分布 $\delta_2=2\delta/15=0.1333\delta$；
线性分布才是 $\delta/6$；位移厚度 $\delta_1=\int_0^{\delta}(1-v_x/v_\infty)\mathrm{d}y=\delta/3$
（抛物线分布），三者不能混用。

\textbf{易错点二：$\tau_w$ 必须取壁面（$y=0$）处的速度梯度。}
抛物线分布在壁面处 $\left(\partial v_x/\partial y\right)_0=2v_\infty/\delta$，
故 $\tau_w=2\mu v_\infty/\delta$；
若想当然地取"平均梯度"$v_\infty/\delta$，则 $\delta=3.87x/\sqrt{\Rey_x}$，比正确值小 29%。

\textbf{易错点三：$C_{\mathrm{Df}}$ 的分母用"一面面积"。}
本题得到 $1.46/\sqrt{\Rey_l}$；若误用两面面积 $2bL$，会得到 $0.73/\sqrt{\Rey_l}$，
与教材布拉修斯解的 $1.328/\sqrt{\Rey_l}$ 相差一倍，容易被被误判为"算错"。
\end{yicuo}

\noindent\textbf{第 8.4 题（孔珑 9-12）\quad 正弦速度分布下层流边界层的厚度与摩擦阻力系数}

\begin{jie}
\textbf{已知：}平板层流边界层内速度分布为正弦曲线
\[ v_x=v_\infty\sin\left(\frac{\pi y}{2\delta}\right) \]
来流速度 $v_\infty$、运动黏度 $\nu$；平板长 $L$、宽 $b$。

\textbf{求：}边界层厚度 $\delta$、摩擦阻力系数 $C_{\mathrm{Df}}$ 与雷诺数 $\Rey$ 的关系式。

\textbf{依据：}$\tau_w=\mu\left(\partial v_x/\partial y\right)_{y=0}$ 与动量积分方程
$\mathrm{d}\delta_2/\mathrm{d}x=\tau_w/(\rho v_\infty^{2})$（与第 8.3 题同）。

\textbf{第 1 步\quad 壁面切应力}
\[ \frac{\partial v_x}{\partial y}=v_\infty\frac{\pi}{2\delta}\cos\left(\frac{\pi y}{2\delta}\right)
\ \Longrightarrow\ \left(\frac{\partial v_x}{\partial y}\right)_{y=0}=\frac{\pi v_\infty}{2\delta} \]
\[ \tau_w=\frac{\pi\mu v_\infty}{2\delta} \qquad (8\text{-}2) \]

\textbf{第 2 步\quad 动量损失厚度}
\[ \delta_2=\delta\int_0^{1}\sin\frac{\pi\eta}{2}\left(1-\sin\frac{\pi\eta}{2}\right)\mathrm{d}\eta
=\delta\left[\frac{2}{\pi}-\frac{1}{2}\right]=\frac{4-\pi}{2\pi}\delta=0.1366\delta \]
（原书取 $0.137\delta$，与精确值 $0.1366\delta$ 一致。）

\textbf{第 3 步\quad 代入动量积分方程}
\[ \frac{4-\pi}{2\pi}\frac{\mathrm{d}\delta}{\mathrm{d}x}
=\frac{\pi\mu v_\infty}{2\delta\rho v_\infty^{2}}=\frac{\pi\nu}{2v_\infty\delta} \]
化简得
\[ \delta\,\mathrm{d}\delta=\frac{\pi^{2}}{4-\pi}\frac{\nu}{v_\infty}\mathrm{d}x \]

\textbf{第 4 步\quad 积分并定积分常数}
\[ \frac{1}{2}\delta^{2}=\frac{\pi^{2}}{4-\pi}\frac{\nu}{v_\infty}x+C \]
边界条件：$x=0$ 处 $\delta=0$，故 $C=0$。
\[ \delta=\sqrt{\frac{2\pi^{2}}{4-\pi}}\sqrt{\frac{\nu x}{v_\infty}}
=4.79\sqrt{\frac{\nu x}{v_\infty}}=\frac{4.79x}{\sqrt{\Rey_x}} \]

\textbf{第 5 步\quad 壁面切应力与雷诺数的关系}
\[ \tau_w=\frac{\pi\mu v_\infty}{2\delta}=\frac{\pi\mu v_\infty}{2\times4.79}\sqrt{\frac{v_\infty}{\nu x}}
=0.328\,\rho v_\infty^{2}\Rey_x^{-1/2} \]

\textbf{第 6 步\quad 一面的总摩擦阻力与摩擦阻力系数}
\[ F_{D\mathrm{f}}=b\int_0^{L}\tau_w\mathrm{d}x
=0.328b\rho v_\infty^{2}\sqrt{\frac{\nu}{v_\infty}}\cdot2\sqrt{L}
=0.656\,bL\rho v_\infty^{2}\Rey_l^{-1/2} \]
\[ C_{\mathrm{Df}}=\frac{F_{D\mathrm{f}}}{\frac{1}{2}\rho v_\infty^{2}bL}=\frac{1.312}{\sqrt{\Rey_l}} \]

\textbf{结果：}$\delta=\dfrac{4.79x}{\sqrt{\Rey_x}}$；
$\tau_w=0.328\rho v_\infty^{2}\Rey_x^{-1/2}$；
$F_{D\mathrm{f}}=0.656\,bL\rho v_\infty^{2}\Rey_l^{-1/2}$；
$C_{\mathrm{Df}}=\dfrac{1.312}{\sqrt{\Rey_l}}$。

\textbf{量纲自检：}$\delta_2$ 为 $\mathrm{m}$；$\tau_w/(\rho v_\infty^{2})$ 为 $\mathrm{m}$ $\checkmark$；
$C_{\mathrm{Df}}$ 无量纲 $\checkmark$。
\end{jie}

\begin{yicuo}
\textbf{易错点一：与抛物线分布的结果对比要记住数量级。}
抛物线分布 $\delta_2=0.1333\delta$、$\delta=5.48x/\sqrt{\Rey_x}$；
正弦分布 $\delta_2=0.1366\delta$、$\delta=4.79x/\sqrt{\Rey_x}$。
两者的 $\delta_2$ 只差 2.5%，但厚度系数相差 14%，
说明"假定速度剖面"的差异会直接传递到边界层厚度上，这也是动量积分方程法的固有误差（约 10%）。

\textbf{易错点二：正弦剖面满足两条基本边界条件。}
$y=0$ 处 $v_x=0$、$y=\delta$ 处 $\partial v_x/\partial y=0$（与外部势流光滑衔接）；
抛物线分布同样满足，而线性分布不满足后者，所以线性分布的结果最不可靠。

\textbf{易错点三：两个积分不要记反。}
$\displaystyle\int_0^{1}\sin\frac{\pi\eta}{2}\mathrm{d}\eta=\frac{2}{\pi}$，
$\displaystyle\int_0^{1}\sin^{2}\frac{\pi\eta}{2}\mathrm{d}\eta=\frac{1}{2}$，
故 $\delta_2=\delta\left(\frac{2}{\pi}-\frac{1}{2}\right)$，不要写成 $\delta\left(1-\frac{2}{\pi}\right)$。
\end{yicuo}

\noindent\textbf{第 8.5 题（孔珑 9-14）\quad 相同雷诺数下水与空气绕流平板的摩擦阻力之比}

\begin{jie}
\textbf{已知：}同一块长 $l$、宽 $b$ 的平板分别放在 $20\ ^{\circ}\mathrm{C}$ 水和 $30\ ^{\circ}\mathrm{C}$ 空气中，
两种情况下平板末端雷诺数 $\Rey_l$ 相同；
$20\ ^{\circ}\mathrm{C}$ 水：$\mu_w=1.005\times10^{-3}\ \mathrm{Pa\cdot s}$、
$\nu_w=1.007\times10^{-6}\ \mathrm{m^2/s}$；
$30\ ^{\circ}\mathrm{C}$ 空气：$\mu_a=1.86\times10^{-5}\ \mathrm{Pa\cdot s}$、
$\nu_a=1.6\times10^{-5}\ \mathrm{m^2/s}$。

\textbf{求：}两种情况下平板摩擦阻力之比（$20\ ^{\circ}\mathrm{C}$ 水的摩擦阻力记为 $F_w$，
$30\ ^{\circ}\mathrm{C}$ 空气的记为 $F_a$）。

\textbf{依据：}层流平板摩擦阻力
\[ F_{D\mathrm{f}}=C_{\mathrm{Df}}\frac{1}{2}\rho v_\infty^{2}bl,\qquad
C_{\mathrm{Df}}=\frac{1.328}{\sqrt{\Rey_l}} \]
$\Rey_l$ 相同时 $C_{\mathrm{Df}}$ 相同；由 $\Rey_l=v_\infty l/\nu$ 得 $v_\infty=\Rey_l\nu/l\propto\nu$。

\textbf{第 1 步\quad 阻力与流体物性的关系}
\[ F_{D\mathrm{f}}\propto\rho v_\infty^{2}\propto\rho\nu^{2}=(\rho\nu)\nu=\mu\nu \]

\textbf{第 2 步\quad 代入两种流体的物性}
\[ \frac{F_w}{F_a}
=\frac{\mu_w\nu_w}{\mu_a\nu_a}
=\frac{1.005\times10^{-3}\times1.007\times10^{-6}}{1.86\times10^{-5}\times1.6\times10^{-5}}
=\frac{1.012\times10^{-9}}{2.976\times10^{-10}}=3.4 \]

\textbf{结果：}相同雷诺数下，$20\ ^{\circ}\mathrm{C}$ 水绕流平板的摩擦阻力约为
$30\ ^{\circ}\mathrm{C}$ 空气的 $\mathbf{3.4}$ 倍。

\textbf{量纲自检：}比值为两个同量纲量之比，无量纲 $\checkmark$；
$\mu\nu=[\mathrm{Pa\cdot s}][\mathrm{m^2/s}]=[\mathrm{kg/(m\cdot s)}][\mathrm{m^2/s}]$，
与 $\rho\nu^{2}=[\mathrm{kg/m^3}][\mathrm{m^2/s^2}]$ 量纲相同 $\checkmark$。
\end{jie}

\begin{yicuo}
\textbf{易错点一：不能只比密度。}
若忽略"雷诺数相同"这个条件直接比密度，会得到 $998/1.165\approx857$，差 250 倍。
关键在：$\Rey_l$ 相同而 $\nu$ 不同时，\textbf{来流速度必须随运动黏度一起变}，$v_\infty\propto\nu$，而阻力含 $v_\infty^{2}$。

\textbf{易错点二：结论是 $F_{D\mathrm{f}}\propto\mu\nu$，两个黏度都要乘。}
由 $F\propto\rho v_\infty^{2}$ 与 $v_\infty\propto\nu$ 得 $F\propto\rho\nu^{2}=\mu\nu$；
只乘动力黏度（或只乘运动黏度）都会得到错误比值。

\textbf{易错点三：物性值要与题给温度对应。}
水取 $20\ ^{\circ}\mathrm{C}$、空气取 $30\ ^{\circ}\mathrm{C}$；
本题之所以比值只有 3.4，是因为水的 $\nu$ 虽小（$1.007\times10^{-6}$），
但水的 $\mu$ 比空气大约 54 倍，两者相乘后只放大到 3.4 倍。
\end{yicuo}

\noindent\textbf{第 8.6 题（孔珑 9-15）\quad 平板混合边界层的摩擦阻力}

\begin{jie}
\textbf{已知：}平板长 $L=6\ \mathrm{m}$、宽 $b=2\ \mathrm{m}$，平行于气流静止安放；
来流速度 $v_\infty=60\ \mathrm{m/s}$；$40\ ^{\circ}\mathrm{C}$ 空气
$\nu=16.9\times10^{-6}\ \mathrm{m^2/s}$、$\rho=1.128\ \mathrm{kg/m^3}$；
层流转捩的临界雷诺数 $\Rey_{x\mathrm{cr}}=10^{6}$。

\textbf{求：}平板的摩擦阻力 $F_{D\mathrm{f}}$。

\textbf{依据：}
① 转捩点位置 $x_{\mathrm{cr}}=\Rey_{x\mathrm{cr}}\nu/v_\infty$，据此判断是否混合边界层；
② 混合边界层（前段层流＋后段湍流）的总摩阻系数按教材式
\[ C_{\mathrm{Df}}=\frac{0.455}{\left(\lg \Rey_l\right)^{2.58}}-\frac{A}{\Rey_l} \]
其中常数 $A$ 随 $\Rey_{x\mathrm{cr}}$ 变化，本题由教材表 9-1 查得 $A=3300$；
③ 板面阻力 $F_{D\mathrm{f}}=C_{\mathrm{Df}}\cdot\dfrac{1}{2}\rho v_\infty^{2}\cdot bL$。

\textbf{第 1 步\quad 求转捩点位置}
\[ x_{\mathrm{cr}}=\frac{\Rey_{x\mathrm{cr}}\nu}{v_\infty}
=\frac{10^{6}\times16.9\times10^{-6}}{60}=0.2817\ \mathrm{m} \]
层流段只有 $0.28\ \mathrm{m}$，其余 $5.72\ \mathrm{m}$ 均为湍流边界层，
是典型的\textbf{混合边界层}，不能按全层流或全湍流处理。

\textbf{第 2 步\quad 求平板末端雷诺数。}
按题给板长 $L=6\ \mathrm{m}$：
\[ \Rey_l=\frac{v_\infty L}{\nu}=\frac{60\times6}{16.9\times10^{-6}}=2.13\times10^{7} \]
原书此式的分子印作 $60\times20$（相当于把板长取成 $20\ \mathrm{m}$），得 $\Rey_l=7.1\times10^{7}$，
与题给 $6\ \mathrm{m}$ 不自洽。下面\textbf{两种口径都算}，并比较结果。

\textbf{第 3 步\quad 求混合边界层摩阻系数。}
口径一（原书，$\Rey_l=7.1\times10^{7}$）：
\[ \lg\left(7.1\times10^{7}\right)=7.851,\qquad (7.851)^{2.58}=203.7 \]
\[ C_{\mathrm{Df}}=\frac{0.455}{203.7}-\frac{3300}{7.1\times10^{7}}
=2.234\times10^{-3}-4.65\times10^{-5}=2.187\times10^{-3} \]
口径二（按题给 $L=6\ \mathrm{m}$，$\Rey_l=2.13\times10^{7}$）：
\[ \lg\left(2.13\times10^{7}\right)=7.328,\qquad (7.328)^{2.58}=170.6 \]
\[ C_{\mathrm{Df}}=\frac{0.455}{170.6}-\frac{3300}{2.13\times10^{7}}
=2.668\times10^{-3}-1.55\times10^{-4}=2.513\times10^{-3} \]

\textbf{第 4 步\quad 求板面摩擦阻力。}
先算动压与面积的乘积（按平板一面 $bL=2\times6=12\ \mathrm{m^2}$）：
\[ \frac{1}{2}\rho v_\infty^{2}bL=\frac{1}{2}\times1.128\times60^{2}\times2\times6=2.4365\times10^{4}\ \mathrm{N} \]
口径一：$F_{D\mathrm{f}}=2.187\times10^{-3}\times2.4365\times10^{4}=53.3\ \mathrm{N}$；
口径二：$F_{D\mathrm{f}}=2.513\times10^{-3}\times2.4365\times10^{4}=61.2\ \mathrm{N}$。

\textbf{结果：}按原书口径（$\Rey_l=7.1\times10^{7}$）$F_{D\mathrm{f}}=53.285\ \mathrm{N}$；
按题给板长复算（$\Rey_l=2.13\times10^{7}$）$F_{D\mathrm{f}}=61.2\ \mathrm{N}$。
若平板两面都迎风，上述数值还应乘以 2。

\textbf{量纲自检：}$\lg \Rey_l$ 为纯数，$C_{\mathrm{Df}}$ 无量纲 $\checkmark$；
$F_{D\mathrm{f}}=[\mathrm{kg/m^3}][\mathrm{m^2/s^2}][\mathrm{m^2}]=\mathrm{kg\cdot m/s^2}=\mathrm{N}$ $\checkmark$。
\end{jie}

\begin{yicuo}
\textbf{易错点一：$\Rey_{x\mathrm{cr}}$ 是局部雷诺数，不是板末端雷诺数。}
局部雷诺数 $\Rey_x=v_\infty x/\nu$ 达到 $10^{6}$ 的位置才是转捩点 $x_{\mathrm{cr}}$。
必须先把 $x_{\mathrm{cr}}$ 算出来（本题 $0.2817\ \mathrm{m}$），
才能判断"层流段＋湍流段"，进而选用混合边界层公式。

\textbf{易错点二：混合边界层公式两项的位置不能颠倒。}
$\dfrac{0.455}{(\lg\Rey_l)^{2.58}}$ 是\textbf{按板末端雷诺数}算的湍流项，
减去第二项 $\dfrac{A}{\Rey_l}$ 才是"前段层流比后段湍流阻力小"的修正；
公式中的 $\Rey$ 一律用板末端的 $\Rey_l$。

\textbf{易错点三：常数 $A$ 随 $\Rey_{x\mathrm{cr}}$ 变化。}
本题 $\Rey_{x\mathrm{cr}}=10^{6}$ 取 $A=3300$；
若 $\Rey_{x\mathrm{cr}}=5\times10^{5}$（真题常用），$A$ 约为 $1700$。
查错表或乱代 $A$ 值会使结果偏差 30% 以上。
\end{yicuo}

\noindent\textbf{第 8.7 题（孔珑 9-16）\quad 用卡门涡街流量计测量管道空气流量}

\begin{jie}
\textbf{已知：}管道直径 $D=500\ \mathrm{mm}=0.5\ \mathrm{m}$，管内为 $30\ ^{\circ}\mathrm{C}$ 空气
（$\nu=15.95\times10^{-6}\ \mathrm{m^2/s}$）；
垂直于管道轴线插入的涡街发生体直径 $d=10\ \mathrm{mm}=0.01\ \mathrm{m}$；
测得漩涡脱落频率 $f=105\ \mathrm{s^{-1}}$；取斯特劳哈尔数 $Sr=0.21$。

\textbf{求：}管道中的体积流量 $q_V$。

\textbf{依据：}卡门涡街的斯特劳哈尔数关系
\[ f=Sr\frac{v}{d}\ \Longrightarrow\ v=\frac{fd}{Sr} \]
（$v$ 为绕过发生体的流速）；流量 $q_V=vA$，管道过流面积 $A=\pi D^{2}/4$。

\textbf{第 1 步\quad 求绕柱流速}
\[ v=\frac{fd}{Sr}=\frac{105\times0.01}{0.21}=5\ \mathrm{m/s} \]

\textbf{第 2 步\quad 验算雷诺数（$Sr=0.21$ 只在大雷诺数下成立）}
\[ \Rey=\frac{vd}{\nu}=\frac{5\times0.01}{15.95\times10^{-6}}=3.13\times10^{3}>10^{3} \]
满足大雷诺数条件，取 $Sr=0.21$ 成立。

\textbf{第 3 步\quad 求管道流量}
\[ q_V=v\cdot\frac{\pi D^{2}}{4}=5\times\frac{\pi\times0.5^{2}}{4}
=5\times0.1963=0.98\ \mathrm{m^3/s} \]

\textbf{结果：}管道中的体积流量 $q_V=0.98\ \mathrm{m^3/s}$。

\textbf{量纲自检：}$fd/Sr=[\mathrm{s^{-1}}][\mathrm{m}]=\mathrm{m/s}$ $\checkmark$；
$q_V=[\mathrm{m/s}][\mathrm{m^2}]=\mathrm{m^3/s}$ $\checkmark$；$\Rey$ 无量纲 $\checkmark$。
\end{jie}

\begin{yicuo}
\textbf{易错点一：$Sr$ 定义式中的特征长度是发生体直径 $d$，算流量却要用管道直径 $D$。}
$d=10\ \mathrm{mm}$、$D=500\ \mathrm{mm}$，两者相差 50 倍；
把管道面积误算成 $\pi d^{2}/4$ 会得到 $0.00039\ \mathrm{m^3/s}$，小 2500 倍。

\textbf{易错点二：先假设后验算。}
$Sr=0.21$ 是大雷诺数（$\Rey>10^{3}$）下圆柱绕流的近似常数，
必须先算出 $\Rey$ 并说明满足条件，才能使用该值。

\textbf{易错点三：涡街流量计直接测的是流速。}
频率算出的 $v$ 是流速，还要乘以管道过流面积才是体积流量（本题即 $5\times0.1963$）。

\textbf{补充：}本题的 $Sr$ 属超出 818 考纲的扩展内容（教材第 8 章只讲分离与尾涡），
仅作了解，不必记忆 $Sr$ 的数值。
\end{yicuo}

\noindent\textbf{第 8.8 题（孔珑 9-17）\quad 求汽车克服空气阻力所消耗的功率}

\begin{jie}
\textbf{已知：}汽车行驶速度 $v=60\ \mathrm{km/h}$；垂直于运动方向的投影面积（迎风面积）
$A=2\ \mathrm{m^2}$；阻力系数 $C_D=0.3$；静止空气温度 $0\ ^{\circ}\mathrm{C}$，
密度 $\rho=1.293\ \mathrm{kg/m^3}$。

\textbf{求：}汽车克服空气阻力所消耗的功率 $P$。

\textbf{依据：}钝体绕流（压差）阻力
\[ F_D=C_D\frac{1}{2}\rho v^{2}A \]
匀速行驶时阻力功率等于阻力与速度之积
\[ P=F_Dv \]

\textbf{第 1 步\quad 统一单位}
\[ v=60\ \mathrm{km/h}=\frac{60\times10^{3}}{3600}=16.67\ \mathrm{m/s} \]

\textbf{第 2 步\quad 求空气阻力}
\[ F_D=0.3\times\frac{1}{2}\times1.293\times16.67^{2}\times2
=0.3\times60.25\times3.6^{2}\times\frac{1}{2}\times\ldots \]
为避免中间取舍误差，按原式计算：
\[ F_D=0.3\times\frac{1}{2}\times1.293\times\left(\frac{60\times10^{3}}{3600}\right)^{2}\times2
=0.3\times0.5\times1.293\times277.8\times2=107.75\ \mathrm{N} \]

\textbf{第 3 步\quad 求功率}
\[ P=F_Dv=107.75\times16.67=1795.8\ \mathrm{W}=1.796\times10^{3}\ \mathrm{W}\approx1.80\ \mathrm{kW} \]

\textbf{结果：}空气阻力 $F_D=107.75\ \mathrm{N}$；
克服空气阻力消耗的功率 $P=1.796\times10^{3}\ \mathrm{W}\approx1.80\ \mathrm{kW}$。

\textbf{量纲自检：}$F_D=[\mathrm{kg/m^3}][\mathrm{m^2/s^2}][\mathrm{m^2}]
=\mathrm{kg\cdot m/s^2}=\mathrm{N}$ $\checkmark$；
$P=[\mathrm{N}][\mathrm{m/s}]=\mathrm{kg\cdot m^2/s^3}=\mathrm{W}$ $\checkmark$。
\end{jie}

\begin{yicuo}
\textbf{易错点一：$60\ \mathrm{km/h}$ 必须先化成 $16.67\ \mathrm{m/s}$。}
功率含速度的三次方（$P=C_D\frac{1}{2}\rho v^{3}A$），
单位没换会差 $(1000/3600)^{3}$ 倍，结果荒谬。

\textbf{易错点二：$A$ 是迎风投影面积，不是车身表面总面积。}
题给"垂直于运动方向的投影面积 $2\ \mathrm{m^2}$"，
若误用车身表面积（几十平方米）会算大十几倍。

\textbf{易错点三：$0\ ^{\circ}\mathrm{C}$ 的空气密度取 $1.293\ \mathrm{kg/m^3}$。}
温度升高密度减小，同车速下阻力也随之减小；
若照搬 $20\ ^{\circ}\mathrm{C}$ 的 $1.205\ \mathrm{kg/m^3}$，阻力会偏小约 7%。
\end{yicuo}

\noindent\textbf{第 8.9 题（孔珑 9-18）\quad 求风洞中垂直吹向圆盘的力}

\begin{jie}
\textbf{已知：}$20\ ^{\circ}\mathrm{C}$ 空气 $\rho=1.205\ \mathrm{kg/m^3}$、
$\nu=15\times10^{-6}\ \mathrm{m^2/s}$；风速 $v=10\ \mathrm{m/s}$，
圆盘直径 $d=50\ \mathrm{cm}=0.5\ \mathrm{m}$，圆盘面垂直于气流方向。

\textbf{求：}作用在圆盘上的力 $F_D$。

\textbf{依据：}先算 $\Rey=vd/\nu$，再由教材图 9-26
（钝体阻力系数与雷诺数的关系曲线）查阻力系数 $C_D$；阻力按
\[ F_D=C_D\frac{1}{2}\rho v^{2}A \]
计算，圆盘迎面投影面积 $A=\pi d^{2}/4$。

\textbf{第 1 步\quad 求雷诺数}
\[ \Rey=\frac{vd}{\nu}=\frac{10\times0.5}{15\times10^{-6}}=3.33\times10^{5} \]

\textbf{第 2 步\quad 查阻力系数}
由图 9-26，在 $\Rey=3.33\times10^{5}$ 处查得 $C_D=1.1$
（钝体绕流在此雷诺数范围内阻力系数变化平缓，可视为常数）。

\textbf{第 3 步\quad 求迎面投影面积}
\[ A=\frac{\pi d^{2}}{4}=\frac{\pi\times0.5^{2}}{4}=0.1963\ \mathrm{m^2} \]

\textbf{第 4 步\quad 求阻力}
\[ F_D=C_D\frac{1}{2}\rho v^{2}A=1.1\times\frac{1}{2}\times1.205\times10^{2}\times0.1963
=1.1\times60.25\times0.1963=13.0\ \mathrm{N} \]

\textbf{结果：}作用在圆盘上的力 $F_D=13.0\ \mathrm{N}$。

\textbf{量纲自检：}$F_D=[\mathrm{kg/m^3}][\mathrm{m^2/s^2}][\mathrm{m^2}]=\mathrm{N}$ $\checkmark$；
$\Rey$ 无量纲 $\checkmark$。
\end{jie}

\begin{yicuo}
\textbf{易错点一：迎风面积按圆面积 $\pi d^{2}/4$ 算。}
圆盘正对气流，投影是直径为 $d$ 的圆，不是 $d^{2}$；
用 $d^{2}$ 会把结果放大 $4/\pi\approx1.27$ 倍。

\textbf{易错点二：这是钝体绕流的压差（形状）阻力，必须查 $C_D$ 曲线。}
不能用层流平板的摩擦阻力系数 $1.328/\sqrt{\Rey_l}$，
后者是"摩擦阻力系数"，与本题的"形状阻力系数"是完全不同的物理量。

\textbf{易错点三：圆球与圆盘的阻力特性不同。}
圆球在 $\Rey\approx3\times10^{5}$ 附近会出现"阻力危机"（$C_D$ 骤降到约 0.2），
而正面迎风的圆盘没有这种骤降，$C_D$ 仍约为 1.1，
两者不能互相套用曲线读数。
\end{yicuo}

\noindent\textbf{第 8.10 题（孔珑 9-19）\quad 由钢球沉降速度求油的黏度并验算雷诺数}

\begin{jie}
\textbf{已知：}透平油密度 $\rho=899.4\ \mathrm{kg/m^3}$；
小钢球直径 $d=3\ \mathrm{mm}=3\times10^{-3}\ \mathrm{m}$、密度 $\rho_s=7791\ \mathrm{kg/m^3}$；
钢球在油中自由降落，测得等速沉降速度 $u=11\ \mathrm{cm/s}=0.11\ \mathrm{m/s}$；
取 $g=9.8\ \mathrm{m/s^2}$。

\textbf{求：}（1）油的动力黏度 $\mu$；（2）验算雷诺数 $\Rey$ 是否在教材式(9-59)的适用范围内。

\textbf{依据：}在斯托克斯区（$\Rey<1$）圆球阻力系数 $C_D=24/\Rey$，
阻力 $F_D=C_D\dfrac{1}{2}\rho u^{2}\dfrac{\pi d^{2}}{4}=3\pi\mu u d$；
钢球重力 $G=\dfrac{1}{6}\pi d^{3}\rho_s g$，浮力 $F_B=\dfrac{1}{6}\pi d^{3}\rho g$；
等速沉降时三力平衡 $F_D+F_B=G$。

\textbf{第 1 步\quad 化简阻力表达式}
\[ F_D=\frac{24\mu}{\rho u d}\times\frac{1}{2}\rho u^{2}\times\frac{\pi d^{2}}{4}=3\pi\mu u d \]

\textbf{第 2 步\quad 列受力平衡式}
\[ F_D+F_B=G\ \Longrightarrow\ 3\pi\mu u d+\frac{1}{6}\pi d^{3}\rho g=\frac{1}{6}\pi d^{3}\rho_s g \]
两边同除以 $\pi d$：
\[ 3\mu u+\frac{1}{6}d^{2}\rho g=\frac{1}{6}d^{2}\rho_s g \]

\textbf{第 3 步\quad 解出黏度}
\[ \mu=\frac{d^{2}\left(\rho_s-\rho\right)g}{18u} \]

\textbf{第 4 步\quad 代入数值}
浮力项：
\[ \frac{1}{6}d^{2}\rho g=\frac{1}{6}\times\left(3\times10^{-3}\right)^{2}\times899.4\times9.8=0.01322\ \mathrm{N} \]
重力项：
\[ \frac{1}{6}d^{2}\rho_s g=\frac{1}{6}\times\left(3\times10^{-3}\right)^{2}\times7791\times9.8=0.11453\ \mathrm{N} \]
故
\[ 3\mu u=0.11453-0.01322=0.10131\ \mathrm{N}
\ \Longrightarrow\ \mu=\frac{0.10131}{3\times0.11}=0.30700\ \mathrm{Pa\cdot s} \]
直接代入导出式复核：
\[ \mu=\frac{\left(3\times10^{-3}\right)^{2}\times(7791-899.4)\times9.8}{18\times0.11}
=\frac{0.60784}{1.98}=0.30700\ \mathrm{Pa\cdot s} \]
与原书给出的 $\mu=0.306989\ \mathrm{Pa\cdot s}$ 一致。

\textbf{第 5 步\quad 验算雷诺数}
\[ \Rey=\frac{\rho u d}{\mu}=\frac{899.4\times0.11\times3\times10^{-3}}{0.30700}
=\frac{0.29680}{0.30700}=0.967<1 \]
$\Rey=0.967<1$，满足教材式(9-59)（斯托克斯区 $C_D=24/\Rey$）的适用范围。

\textbf{结果：}透平油的动力黏度 $\mu=0.30700\ \mathrm{Pa\cdot s}$（原书 $0.306989\ \mathrm{Pa\cdot s}$）；
$\Rey=0.967<1$，斯托克斯公式适用。

\textbf{量纲自检：}$\dfrac{d^{2}(\rho_s-\rho)g}{18u}
=\dfrac{[\mathrm{m^2}][\mathrm{kg/m^3}][\mathrm{m/s^2}]}{[\mathrm{m/s}]}
=\mathrm{kg/(m\cdot s)}=\mathrm{Pa\cdot s}$，量纲正确 $\checkmark$。
\end{jie}

\begin{yicuo}
\textbf{易错点一：等速沉降是"三力平衡"，浮力不能漏。}
本题浮力 $0.01322\ \mathrm{N}$ 约占重力的 11.5%，
漏掉它会使算得的 $\mu$ 偏大约 13%（$0.347\ \mathrm{Pa\cdot s}$）。

\textbf{易错点二：$C_D=24/\Rey$ 代入后阻力与速度成一次关系。}
$F_D=3\pi\mu u d$ 是斯托克斯公式，
与湍流区的"阻力正比于 $u^{2}$"完全不同；
若误用 $C_D$ 为常数的牛顿区公式，就把可解的方程变成了非线性的。

\textbf{易错点三：算出黏度后必须回代验算 $\Rey<1$。}
$C_D=24/\Rey$ 只在 $\Rey<1$ 时成立，本题 $\Rey=0.967$ 已很接近 1，
说明钢球的直径和油的黏度都恰好使得该式勉强适用；
若换成更大的钢球或更稀的油，就必须改用其它分区的 $C_D$ 重新试算。
\end{yicuo}

\noindent\textbf{第 8.11 题（孔珑 9-20）\quad 求使料层颗粒悬浮的最小风速}

\begin{jie}
\textbf{已知：}采暖沸腾炉料层温度 $1000\ ^{\circ}\mathrm{C}$；
烟气运动黏度 $\nu=1.67\times10^{-6}\ \mathrm{m^2/s}$；
标准状态烟气密度 $\rho_0=1.34\ \mathrm{kg/m^3}$；
颗粒密度 $\rho_p=2294\ \mathrm{kg/m^3}$、平均粒径 $d=1.68\ \mathrm{mm}=1.68\times10^{-3}\ \mathrm{m}$；
取 $g=9.8\ \mathrm{m/s^2}$。

\textbf{求：}使颗粒处于悬浮状态所需的最小风速 $u$。

\textbf{依据：}烟气自下而上绕流颗粒，颗粒受阻力
$F_D=C_D\dfrac{1}{2}\rho u^{2}\dfrac{\pi d^{2}}{4}$、浮力 $F_B=\dfrac{1}{6}\pi d^{3}\rho g$、
重力 $G=\dfrac{1}{6}\pi d^{3}\rho_p g$；
颗粒刚好悬浮（即将被吹起、失去静止支撑）时
\[ F_D+F_B=G \]
阻力系数按雷诺数分区取值：$\Rey<1$ 取 $C_D=24/\Rey$；
$1000<\Rey<2\times10^{5}$ 取 $C_D=0.45$（教材 8.7 节）。

\textbf{第 1 步\quad 求 $1000\ ^{\circ}\mathrm{C}$ 烟气的密度}
烟气按等压理想气体由标准状态折算：
\[ \rho=\rho_0\frac{273}{273+t}=1.34\times\frac{273}{273+1000}
=1.34\times\frac{273}{1273}=0.2874\ \mathrm{kg/m^3} \]

\textbf{第 2 步\quad 化简悬浮条件}
\[ C_D\frac{1}{2}\rho u^{2}\frac{\pi d^{2}}{4}+\frac{1}{6}\pi d^{3}\rho g
=\frac{1}{6}\pi d^{3}\rho_p g \]
两边同除以 $\dfrac{1}{4}\pi d^{2}$（或同除以 $\frac{1}{6}\pi d^{2}$ 亦可），得
\[ \frac{1}{2}C_D\rho u^{2}=\frac{2}{3}d\left(\rho_p-\rho\right)g
\ \Longrightarrow\ u=\sqrt{\frac{4d\left(\rho_p-\rho\right)g}{3C_D\rho}} \]

\textbf{第 3 步\quad 先假设 $\Rey<1$（斯托克斯区），取 $C_D=24/\Rey$}
此时上式化为
\[ u=\frac{d^{2}\left(\rho_p-\rho\right)g}{18\mu},\qquad
\mu=\rho\nu=0.2874\times1.67\times10^{-6}=4.80\times10^{-7}\ \mathrm{Pa\cdot s} \]
\[ u=\frac{\left(1.68\times10^{-3}\right)^{2}\times(2294-0.2874)\times9.8}
{18\times4.80\times10^{-7}}=\frac{0.06344}{8.64\times10^{-6}}=7.34\times10^{3}\ \mathrm{m/s} \]

\textbf{第 4 步\quad 验算雷诺数，发现假设不成立}
\[ \Rey=\frac{ud}{\nu}=\frac{7.34\times10^{3}\times1.68\times10^{-3}}{1.67\times10^{-6}}
=7.39\times10^{6}\gg1 \]
与 $\Rey<1$ 矛盾，斯托克斯区不适用。

\textbf{第 5 步\quad 改设 $1000<\Rey<2\times10^{5}$，取 $C_D=0.45$}
\[ u=\sqrt{\frac{4\times1.68\times10^{-3}\times(2294-0.2874)\times9.8}
{3\times0.45\times0.2874}}=\sqrt{\frac{151.1}{0.3880}}=\sqrt{389.4}=19.7\ \mathrm{m/s} \]

\textbf{第 6 步\quad 再验算}
\[ \Rey=\frac{ud}{\nu}=\frac{19.73\times1.68\times10^{-3}}{1.67\times10^{-6}}=1.98\times10^{4} \]
$1000<\Rey=1.98\times10^{4}<2\times10^{5}$，落在 $C_D=0.45$ 的分区范围内，假设成立。

\textbf{结果：}通过料层所需最小风速 $u=19.7\ \mathrm{m/s}$（约 $20\ \mathrm{m/s}$），
此时颗粒所受到的阻力与浮力之和恰好等于其重力，颗粒处于悬浮状态。

\textbf{量纲自检：}$\sqrt{\dfrac{[\mathrm{m}][\mathrm{kg/m^3}][\mathrm{m/s^2}]}{[\mathrm{kg/m^3}]}}
=\sqrt{\mathrm{m^2/s^2}}=\mathrm{m/s}$，量纲正确 $\checkmark$；$\Rey$ 无量纲 $\checkmark$。
\end{jie}

\begin{yicuo}
\textbf{易错点一：$1000\ ^{\circ}\mathrm{C}$ 烟气的密度必须由状态方程折算。}
高温下烟气密度只有标准状态的 $0.2874/1.34\approx21\%$。
若直接取 $1.34\ \mathrm{kg/m^3}$，由 $u\propto1/\sqrt{\rho}$ 会把风速算小
$\sqrt{1.34/0.2874}=2.16$ 倍（得 9.1 m/s），严重偏低。

\textbf{易错点二：悬浮条件是"阻力＋浮力＝重力"。}
$\rho_p-\rho=2294-0.2874$，浮力项虽只占重力的 0.013%，
但公式中必须写成 $(\rho_p-\rho)$，不能写成 $\rho_p$。

\textbf{易错点三：$C_D$ 必须"先假设、算完再回代验算"。}
本题如果一上来就取 $C_D=0.45$，虽然碰巧得到正确结果，
但完整解答必须给出两个分区的试算过程与验算结论（真题评分按步给分）。
另外，题给的是\textbf{运动黏度} $\nu$，用它算斯托克斯区公式时要先换成
$\mu=\rho\nu$，不要直接把 $\nu$ 当 $\mu$ 用。
\end{yicuo}
